\section{Résultats de l'étape 3}

\subsection{Campagne de mesures}

La campagne officielle a été menée sur \texttt{data/galaxy\_5000}, avec \(dt = 0.0015\), une grille \(15 \times 15 \times 1\), un pas de warmup et 30 pas mesurés. Les couples \((\text{workers}, \text{threads})\) retenus respectent la contrainte \(\text{workers} \times \text{threads} \leq 16\) sur cette machine.

Pour les cas à 9 processus MPI au total, il a fallu ajouter l'option \texttt{--use-hwthread-cpus} à \texttt{mpirun}, car OpenMPI compte par défaut les slots sur les seuls cœurs physiques.

\subsection{Table 2D des temps bout-en-bout}

\begin{table}[H]
    \centering
    \begin{tabular}{c c c c}
        \toprule
        \textbf{Workers MPI} & \textbf{1 thread} & \textbf{2 threads} & \textbf{4 threads} \\
        \midrule
        4 & 125.816 & 72.342 & 52.114 \\
        6 & 163.798 & 100.596 & N/A \\
        8 & 139.213 & Instable & N/A \\
        \bottomrule
    \end{tabular}
    \caption{Temps bout-en-bout moyens de l'étape 3, en ms}
    \label{tab:stage3_matrix}
\end{table}

Dans cette table, \texttt{N/A} signifie qu'une configuration n'a pas été mesurée parce qu'elle dépasserait la limite \(\text{workers} \times \text{threads} \leq 16\). Le terme \og Instable \fg{} désigne au contraire une configuration réellement testée, mais écartée car elle produit des valeurs non finies au cours du calcul.

\subsection{Détails sur les configurations stables}

\begin{table}[H]
    \centering
    \resizebox{\linewidth}{!}{%
    \begin{tabular}{c c c c c c c c}
        \toprule
        \textbf{Workers} & \textbf{Threads} & \textbf{end-to-end (ms)} & \textbf{worker mean (ms)} & \textbf{worker max (ms)} & \textbf{Migration (ms)} & \textbf{Halo (ms)} & \textbf{Réduction (ms)} \\
        \midrule
        4 & 1 & 125.816 & 101.738 & 125.148 & 26.652 & 0.493 & 0.255 \\
        4 & 2 & 72.342 & 58.431 & 71.717 & 12.227 & 0.488 & 0.259 \\
        4 & 4 & 52.114 & 43.161 & 51.457 & 7.620 & 0.537 & 0.276 \\
        6 & 1 & 163.798 & 108.963 & 162.997 & 52.802 & 0.596 & 0.325 \\
        6 & 2 & 100.596 & 68.752 & 99.907 & 32.172 & 0.627 & 0.347 \\
        8 & 1 & 139.213 & 95.797 & 138.461 & 47.931 & 0.658 & 0.336 \\
        \bottomrule
    \end{tabular}%
    }
    \caption{Mesures détaillées des configurations stables de l'étape 3}
    \label{tab:stage3_details}
\end{table}

\subsection{Graphiques}

\begin{figure}[H]
    \centering
    \begin{subfigure}[t]{0.48\linewidth}
        \centering
        \begin{tikzpicture}
\begin{axis}[
    width=\linewidth,
    height=6.5cm,
    xlabel={Nombre de workers MPI},
    ylabel={Speed-up relatif à 4 workers},
    xmin=4,
    xmax=8,
    ymin=0,
    xtick={4,6,8},
    legend style={at={(0.5,-0.2)}, anchor=north, legend columns=3},
    grid=both,
    minor y tick num=1,
    axis line style={draw=black!60},
]
\addplot[color=enstaBleuFonce, mark=* , line width=1.1pt] coordinates {
        (4, 1.000)
        (6, 0.768)
        (8, 0.904)
};
\addlegendentry{1 thread}

\addplot[color=enstaBleuClair, mark=square* , line width=1.1pt] coordinates {
        (4, 1.000)
        (6, 0.719)
};
\addlegendentry{2 threads}

\addplot[color=black!70, mark=triangle* , line width=1.1pt] coordinates {
        (4, 1.000)
};
\addlegendentry{4 threads}
\end{axis}
\end{tikzpicture}

        \caption{Speed-up relatif à 4 workers}
        \label{fig:stage3_speedup}
    \end{subfigure}
    \hfill
    \begin{subfigure}[t]{0.48\linewidth}
        \centering
        \begin{tikzpicture}
\begin{axis}[
    width=\linewidth,
    height=6.5cm,
    xlabel={Nombre de threads Numba},
    ylabel={Temps bout-en-bout (ms)},
    xmin=1,
    xmax=4,
    ymin=0,
    xtick={1,2,4},
    legend style={at={(0.5,-0.2)}, anchor=north, legend columns=2},
    grid=both,
    minor y tick num=1,
    axis line style={draw=black!60},
]
\addplot[color=black!70, mark=square* , line width=1.1pt] coordinates {
    (1, 246.919)
    (2, 138.853)
    (4, 79.767)
};
\addlegendentry{Étape 02}

\addplot[color=enstaBleuFonce, mark=* , line width=1.1pt] coordinates {
    (1, 125.816)
    (2, 72.342)
    (4, 52.114)
};
\addlegendentry{Étape 03 avec 4 workers}
\end{axis}
\end{tikzpicture}

        \caption{Étape 2 vs étape 3 à nombre de threads égal}
        \label{fig:stage2_vs_stage3}
    \end{subfigure}
    \caption{Lecture globale des performances de l'étape 3}
    \label{fig:stage3_global}
\end{figure}

Le résultat principal est simple : l'étape 3 devient plus rapide que l'étape 2 à nombre de threads égal dès que l'on utilise 4 workers. Par exemple, à 4 threads, le temps bout-en-bout passe de \SI{79.767}{ms} à l'étape 2 à \SI{52.114}{ms} à l'étape 3.

Le gain reste net sur les trois configurations stables à 4 workers : environ \(\times 1.96\) à 1 thread, \(\times 1.92\) à 2 threads et \(\times 1.53\) à 4 threads par rapport à l'étape 2.

En revanche, augmenter encore le nombre de workers ne donne pas ici un meilleur speed-up. Les configurations à 6 et 8 workers deviennent moins bonnes que celle à 4 workers, ce qui montre que la décomposition de domaine ne suffit pas à elle seule : le coût de migration et surtout le déséquilibre entre workers prennent alors le dessus.
